% Cote 119: offprints, reports, courses, CV and letters (typed and printed part): archive cote 119, site page 263 (PDF page 264).
% typescript. Draft compared: none.
% AI-transcribed by Claude Opus 5.5 (claude-opus-5-5, adaptive thinking, effort high), Message Batches API, batch r119_full (msgbatch_01DwZNm5Y7YsKBWyAN4sFs8w), 2026-11-13; not yet checked by eye.
\scanpage{119/263}
\mspage{10}

B. \uline{Demain}

Nous esquissons ici quelques-unes des propositions qui ont été avancées, et
pour la plupart discutées, concernant la structure du département de mathé-%
matique, et le système de formation, de qualification et de recrutement des
enseignants et chercheurs en mathématique.

1. Le département serait partagé en deux sections, savoir l'\uline{Institut de
Mathématique}, groupant les mathématiciens à vocation principale de recherche,
et l'\uline{Ecole de Mathématique}, groupant les mathématiciens à vocation princi-%
pale d'enseignement. Le rôle du premier serait, en plus de la recherche,
l'enseignement mathématique à un niveau théorique élevé, plus particulièrement
l'enseignement post-licence, destiné surtout aux futurs mathématiciens, et
la formation des chercheurs. Le rôle de l'Ecole de Mathématique serait l'en-%
seignement mathématique en général, et plus particulièrement l'enseignement
pré-licence à l'usage de futurs usagers non mathématiciens. Les décisions
essentielles (organisation de l'enseignement, recrutement etc.) seraient
prises par une assemblée commune, l'\uline{Assemblée du département de Mathématique},
formée des maîtres de conférences et professeurs des deux sections, de
représentants des étudiants et de représentants des assistants et maîtres-%
assistants. Entre les deux sections du département, les relations seraient
très étroites, notamment par la possibilité pour un membre de l'Ecole de
Mathématique de se faire détacher une année sur cinq à l'Institut de Mathé-%
matique avec dispense d'enseignement, et la possibilité pour un membre de
l'Institut qui le désirerait d'être détaché, avec l'accord de ses collègues
de l'Ecole de Mathématique, pendant une ou plusieurs années à cette dernière.
Les membres de l'une et l'autre section seraient considérés comme exerçant
des métiers essentiellement distincts, quoique étroitement liés,- ce qui
n'empêcherait pas un membre de l'une des sections ayant les qualifications
voulues de demander à changer de section à titre permanent, la décision
incombant alors à l'Assemblée du département de Mathématique, au même titre
que toute autre nomination.

\margin{263}
